\begin{center}
\LARGE
On an Algebraic Characterization of Default Reasoning 
\end{center}

\begin{center}
\Large
Helmut Thiele
\end{center}

\noindent
Date:  July 17th (Wednesday)\\
Time: 16:30-17:30\\

\begin{center}
\large
Abstract
\end{center}

This talk deals with non-monotonic logical consequence operators
defined by default reasoning.  The aim is to consider and investigate
such operators within the general framework of the theory of closure
operators, in particular, topological and algebraic closure operators
used in topology, algebra, logic and metamathematics.The main result
of the paper is an appropriate generalization of the well-known
theorem of {\it G. Birkhoff}, {\it P. Hall} and {\it J. Schmidt},
which gives an algebraic characterization of compact closure operators
by universal algebras and also by classical, "monotonic", deductive
systems.  For this p[urpose we introduce a new type of algebraic
strusture generalizing the classical notion of universal algebra
(introduced by {\it G. Birkhoff}) which we shall call {\it regulated
  algebra}.  Furthermore, generalizing the concept of a classical
monotonic deductive system which is closely asociated with the theorem
of {\it G.  Birkhoff}, {\it P. Hall} and {\it J. Schmidt}, we
introduce {\it Regulated} {\it Deductive} {\it Systems} which include,
for instance {\it Reiter}'s default theories as special cases.  As we shall
see in the sequel, many of the classical results in non-monotonic
reasoning can be derived on the basis of this general axiomatic
approach. Furthermore, this approach gives a deeper mathematical
understanding of what default reasoning is and provides the starting
point for further inverstigations of regulated algebras and regulated
deductive systems.

